% ===== UNIT 12: Stream 1: loss function + equations (3.2.2) =====
% Status: REWRITTEN  (change to REWRITTEN / REVIEWED when done; see CHECKLIST.md)
% Rules: style rewrite only. Keep every \cite, \ref, \label, number, equation and table value as is. Change the wording, not the content.
\subsubsection{Loss function}
Training uses a composite loss. The classification loss is cross-entropy
with label smoothing ($\varepsilon = 0.1$):
\begin{equation}
  \mathcal{L}_{\mathrm{CE\text{-}LS}} =
    -\sum_{c=1}^{C} y_c^{\mathrm{smooth}} \log \hat{p}_c, \;\;
    y_c^{\mathrm{smooth}} = (1{-}\varepsilon)y_c + \varepsilon/C
  \label{eq:cels}
\end{equation}
where $C = 576$ training identities. The recipe of the released headline checkpoint
(A5$\alpha$) adds two terms to this loss. One is a soft-margin triplet loss
(margin $m{=}0.3$). The other is a delayed centre loss ($\lambda{=}5\times10^{-4}$),
enabled after epoch~30. We also use layer-wise learning-rate decay of 0.75:
\begin{equation}
  \mathcal{L}_{\mathrm{A5\alpha}} =
    \mathcal{L}_{\mathrm{CE\text{-}LS}}
    + \mathcal{L}_{\mathrm{tri}}(m{=}0.3)
    + \lambda\,\mathcal{L}_{\mathrm{centre}}
  \label{eq:a5alpha}
\end{equation}
We also evaluate an alternative, a supervised-contrastive variant (A5$\beta$).
It replaces the triplet and centre terms with Supervised Contrastive Loss
(SupCon)~\cite{khosla2020supcon} ($\tau = 0.07$, LLRD 0.65):
\begin{equation}
  \mathcal{L}_{\mathrm{SupCon}} =
    \sum_{i \in \mathcal{B}}
      \frac{-1}{|\mathcal{P}(i)|}
      \sum_{p \in \mathcal{P}(i)}
        \log
        \frac{\exp(\mathbf{z}_i \cdot \mathbf{z}_p / \tau)}
             {\sum_{a \in \mathcal{A}(i)} \exp(\mathbf{z}_i \cdot \mathbf{z}_a / \tau)}
  \label{eq:supcon}
\end{equation}
where $\mathcal{P}(i)$ is the set of positive pairs and $\mathcal{A}(i)$ is all
other mini-batch samples. This gives
$\mathcal{L}_{\mathrm{A5\beta}} = \mathcal{L}_{\mathrm{CE\text{-}LS}} +
\mathcal{L}_{\mathrm{SupCon}}$. A5$\alpha$ and A5$\beta$ are separate ablation rows.
In single from-scratch ablation runs, A5$\alpha$ gets 89.06\% mAP and A5$\beta$ gets 88.57\% mAP
under stream-1-only AQE$+$rerank evaluation (the released A5$\alpha$ checkpoint scores 89.63\% at this setting). The difference is 0.49\,pp, with A5$\alpha$ higher.
This is well inside the $\pm 0.55$\,pp three-seed standard deviation of
A5$\alpha$ (88.75 $\pm$ 0.55\,pp mAP, see Table~\ref{tab:ablation} and
Section~\ref{sec:seed_variance}). So the two recipes are statistically equivalent.
We keep A5$\alpha$ as the headline Stream~1 recipe, because the released checkpoint was trained with it.

CircleLoss made training unstable in our full-recipe runs. The training loss diverged,
so we leave it out of the final recipe. To find the cause, we ran a controlled 400-step
diagnostic with a minimal head (plain CLIP-initialised ViT-B $+$ BNNeck $+$ CE-LS $+$ metric loss,
$P{\times}K=12{\times}8$). We tested five conditions. These were CircleLoss at
$\gamma\in\{32,64\}$ under fp16 and fp32, an ImageNet-initialised control, and
a triplet$+$centre control. None of the CircleLoss conditions diverged
in this minimal setting. The fp16 and fp32 runs differed only in descent speed.
CLIP and ImageNet initialisation behaved similarly. So neither precision nor
initialisation is a simple explanation. The instability we saw is specific to the
full recipe (SIE/JPM heads, layer-wise learning-rate decay and the longer schedule).
It does not come from CircleLoss on CLIP features alone. We did not isolate the exact cause.
CircleLoss stays excluded because of what we observed in the full setting.
The probe script and loss curves are released with the code.
